% =====================================================================
%  PHAN KET QUA — dong gop CHINH cua bai.
%
%  ⚠ KHONG GO SO NAO VAO DAY. Moi con so den tu figures/out/numbers.tex,
%  do figures/make_figures.py sinh ra tu analysis/model.py va results/*.json.
%  Neu can mot so chua co macro thi THEM MACRO, dung go tay.
%
%  Chay `bash figures/build.sh` truoc khi dich ban thao.
% =====================================================================

\section{Measurement Setup}
\label{sec:setup}

All measurements run on a single machine (Ubuntu~24.04) so that every number in this letter
comes from one environment. Two quantities are measured, and neither requires a post-quantum
library: the argument is about message \emph{size}, so credentials are padded or substituted to
post-quantum dimensions rather than computed.

\textbf{Exchanges.} A relay placed between a CoAP client and server counts UDP datagrams in
each direction. Counting datagrams on the wire, rather than invocations inside the server,
matters: the implementation used here caches an assembled body and slices it, so a handler is
called once regardless of how many blocks cross the link.

\textbf{Handshake completion.} A DTLS client and server are run at a fixed MTU with
certificates of increasing size, and a handshake is recorded as complete only on evidence that
a successful run alone can produce, never on the absence of an error. Every implementation is
first exercised against each credential at a wide MTU; a pair that cannot complete there is
excluded rather than swept, since its cells would report a limit of the library rather than of
the link. This excluded one pair, discussed in Section~\ref{sec:threats}.

The block size is fixed at \numCoapBlock~B, the default in both Contiki-NG and RIOT, and the
constrained MTU is \numFramePayload~B, the payload of an IEEE~802.15.4 frame after the MAC
header and AES-CCM$^*$.

\section{Results}
\label{sec:results}

\subsection{The exchange count model agrees with an independent implementation}

Table~\ref{tab:validation} compares the exchange counts predicted by the block-wise model
against datagrams measured with \numMOneImpl, an implementation independent of the model:
\numMOneMatch/\numMOneTotal{} cases agree exactly. This validates the \emph{counting}, and
nothing more. It says nothing about latency or energy on a real link, neither of which is
claimed here.

With that model, a post-quantum EDHOC handshake at ML-KEM-768 and ML-DSA-65 takes
\numExchTotalPQ{} exchanges: \numExchOnePQ{} for message\_1 alone, whose \numMsgOnePQ~B do not
fit in a \numFramePayload~B frame. Fig.~\ref{fig:exchanges} shows the growth.

\subsection{The DTLS comparator, measured rather than cited}

DTLS fragments at its own handshake layer and sends each flight back-to-back, so its
round-trip count should not grow with size \cite{rfc9147}. We measure this rather than assume
it, at the same MTU and with the same relay. Certificate \emph{chains} are used to enlarge the
handshake without enlarging the key, since the library rejects keys beyond its configured
maximum integer size (Section~\ref{sec:threats}).

Across \numDtlsMinFrag{} to \numDtlsMaxFrag{} fragments, \numDtlsImpl{} completes every
handshake, and the datagram count grows from \numDtlsDatagramLo{} to \numDtlsDatagramHi{}
($\times\numDtlsDatagramGrowth$) while the number of direction changes stays at
\numDtlsMeasTurns{} throughout. The flight structure is therefore not an idealisation: at
\numDtlsMaxFrag{} fragments, more than twice what a post-quantum handshake requires, the
round-trip count has not moved.

Two qualifications matter. The measurement is DTLS~\numDtlsMeasVersion, because none of the
three implementations offers DTLS~1.3 over datagrams: GnuTLS lists protocol versions up to
DTLS~1.2 only, mbedTLS accepts \texttt{tls12}, \texttt{dtls12} and \texttt{tls13} but no
\texttt{dtls13}, and OpenSSL's \texttt{s\_server} has no DTLS~1.3 option. Four years after
RFC~9147, the two-round-trip figure it specifies is not obtainable from any of them. And the
constant here is \numDtlsMeasTurns{} direction changes rather than the \numDtlsRT{} round
trips of DTLS~1.3, since DTLS~1.2 carries more flights; what is measured is the invariance,
not the constant's value.

The threshold, not the ratio, is what matters. Classically an EDHOC message fits inside one
frame and block-wise never engages, so both protocols take \numDtlsRT{} round trips and
EDHOC's advantage is expressed purely in bytes. The ML-KEM-512 encapsulation key alone is
\numKemFiveOneTwo~B, which is $\numKemFrameRatio\times$ that frame payload, so block-wise engages in
message\_1 at every NIST parameter set, whatever authentication is chosen. This is a property
of transmitting a KEM public key, not of a particular design.

\subsection{Three implementations, three different limits}

Table~\ref{tab:impls} and Fig.~\ref{fig:impls} report whether each implementation completes a
DTLS handshake as MTU and certificate size vary. The three do not merely differ in degree;
they are bounded by different quantities.

\textbf{OpenSSL is bounded by MTU.} It fails at every MTU at or below \numOpensslMtuBad~B and
completes at every MTU at or above \numOpensslMtuOk~B, for all three certificate sizes. The
boundary sits at OpenSSL's \texttt{DTLS1\_MIN\_MTU} of 256~B. Fragment count explains nothing
here: it fails with \numOpensslFewFragFail{} fragments at MTU~\numOpensslMtuBad{} and completes with \numOpensslManyFragOk{} at MTU~\numOpensslMtuOk{}. At the
\numFramePayload~B payload of an 802.15.4 frame it completes \numOpensslAtFrame{} cases, which
is to say none. \textbf{This holds for classical credentials too}; it is not a post-quantum
result.

\textbf{GnuTLS is bounded by fragment count.} It completes at MTU~\numFramePayload, so no MTU
floor applies, but it completes up to \numGnutlsFragOk{} fragments and fails from
\numGnutlsFragBad{}. The failures are not an artefact of an impatient
harness: given an allowance far beyond the probe timeout, the client still
abandons the handshake on its own retransmission schedule, well inside that allowance. At
\numFramePayload~B, \numGnutlsFragOk{} fragments cap a handshake message at about
\numGnutlsCapKB~kB, which post-quantum credentials exceed.

\textbf{mbedTLS shows neither limit} in the range swept, completing
\numMbedtlsOk/\numMbedtlsTotal{} cases including \numMbedtlsAtFrame{} at
MTU~\numFramePayload{} with \numMbedtlsMaxFrag{} fragments.

This ordering is worth stating plainly. Of the three, mbedTLS is the one actually deployed on
constrained devices, and it is the one that handles the constrained link; the two
server-side implementations do not. The expectation that motivated this measurement was the
opposite.

\subsection{No block size recovers the round trips}

The obvious objection to Section~\ref{sec:results} is that the block size is a tunable
parameter. Fig.~\ref{fig:blocksize} sweeps the full range RFC~7959 permits. Larger blocks do
reduce the exchange count, but each block then exceeds the frame and is fragmented by the
adaptation layer, where a lost fragment costs the whole block~\cite{rfc4944,rfc6282}. RFC~9177~\cite{rfc9177} defines
alternative block options, and dismissing them because no constrained stack implements them
would state a structural conclusion on a contingent fact. We therefore model them.

Q-Block with non-confirmable messages sends blocks serially ``without having to wait for a
response or next request'', acknowledging them in sets of \numQBlockMaxPayloads{}, the default
value of \texttt{MAX\_PAYLOADS}~\cite{rfc9177}. The exchange count becomes the number of such
sets rather than the number of blocks: \numQBlockExch{} instead of \numExchTotalPQ{}, a
reduction of $\numQBlockGain\times$. So the option does recover most of what lock-step costs.

It does not close the gap. \numQBlockExch{} exchanges remain against the
\numDtlsMeasTurns{} direction changes measured for DTLS, and the recovery is unavailable in
practice: neither Contiki-NG nor RIOT implements Q-Block, libcoap gates it behind a flag that
both peers must set, and neither RFC~9528 nor RFC~9668 references RFC~9177. The conclusion this
letter draws is therefore about EDHOC as profiled in RFC~9668, not about every conceivable
CoAP transport for it. The optimum lies at the
boundary of the permitted range rather than inside it, at \numBestBlock~B, and even there the
count is \numBestExch{} exchanges against \numDtlsRT{} DTLS flights, a factor of
$\numBestRatio\times$.

Two further points close the objection. The optimum is not reachable on RIOT, whose
\texttt{CONFIG\_NANOCOAP\_BLOCK\_SIZE\_MAX} is a compile-time ceiling at \numCoapBlock~B; and
the latency-optimal and energy-optimal block sizes are not the same, so no single setting is
best on both axes.
